\documentclass[11pt]{article}
\usepackage[margin=1in]{geometry}
\usepackage{amsmath,amssymb,amsthm,mathtools}
\usepackage{graphicx}
\usepackage{microtype}
\usepackage[colorlinks=true,linkcolor=blue,urlcolor=blue,citecolor=blue]{hyperref}

\newtheorem{theorem}{Theorem}
\newtheorem{corollary}[theorem]{Corollary}
\newcommand{\tr}{\operatorname{tr}}

\title{A superfactorial obstruction to the proposed lower bound for $B_3(2m+1)$}
\author{Candidate disproof of inequality (27) in arXiv:1210.0835}
\date{27 August 2026}

\begin{document}
\maketitle

\noindent\fbox{\parbox{0.94\linewidth}{\textbf{Status: candidate
counterexample/disproof, likely valid.}  The lower bound proposed in
Question~6 cannot hold: the entire absolute path sum is at most
$\binom{m+5}{3}/(4^{m+4}(m!)^2)$, which is superfactorially smaller than
$C^m/m!$ for every fixed $C>0$.}}

\section{Source question}

Djakov and Mityagin study weighted walks arising in the spectral analysis of
Hill operators \cite{DM}.  For steps $F=\{-2,+4\}$, $n=2m+1$, and exactly
three negative steps, their Question~6 on PDF page~5 asks for an asymptotic
formula for $B_3(2m+1)$, ``or at least'' a positive absolute constant $C$
such that
\begin{equation}\label{eq:source-lower}
 |B_3(2m+1)|\geq \frac{C^m}{m!}.
\end{equation}
They note that \eqref{eq:source-lower} would imply a non-basis result for an
associated antiperiodic Hill operator.

\begin{figure}[ht]
\centering
\includegraphics[width=\linewidth]{figures/open_problem_crop.png}
\caption{Question~6 and inequalities (27)--(28) on PDF page~5 of
arXiv:1210.0835v1.}
\end{figure}

\section{Definitions}

An admissible walk $x$ starts at $-n$, ends at $n$, uses steps from
$\{-2,+4\}$, and has no internal vertex equal to either endpoint.  If its
internal vertices are $j(1),\ldots,j(\nu)$, its unweighted path factor is
\[
 h_1(x)=\prod_{t=1}^{\nu}\frac{1}{n^2-j(t)^2}.
\]
Let $X_n^+(3)$ denote the source's collection of such walks with exactly
three $-2$ steps.  Then
\[
 B_3(2m+1)=\sum_{x\in X_n^+(3)}h_1(x),\qquad n=2m+1.
\]
The step-balance equation forces every such walk to have $m+2$ positive
steps.  Hence it has $m+5$ total steps and $m+4$ internal vertices.

\section{Proof intuition}

No cancellation estimate is needed.  Divide positions by two and inspect
only the vertices immediately after the first $m+1$ positive steps.  At the
$p$-th selected vertex, at most three negative steps have occurred, so its
distance factor from the left endpoint grows at least like $p$ and its
distance factor from the right endpoint grows, in reverse order, at least
like $m-p$.  Their products contribute two factorials to every path
denominator.  The number of possible placements of the three negative steps
is only cubic in $m$.  Therefore even the absolute sum is of order at most an
exponential times $(m!)^{-2}$, ruling out the proposed $(m!)^{-1}$ scale.

\section{The obstruction}

\begin{theorem}\label{thm:bound}
For every integer $m\geq1$,
\begin{equation}\label{eq:main-bound}
 |B_3(2m+1)|
 \leq \sum_{x\in X_n^+(3)}|h_1(x)|
 \leq \frac{\binom{m+5}{3}}{4^{m+4}(m!)^2}.
\end{equation}
\end{theorem}

\begin{proof}
Put $n=2m+1$.  Normalize the positions by writing
\[
 j(t)=-n+2r_t.
\]
Thus $r_0=0$, $r_{m+5}=n$, and the normalized steps are $-1$ and $+2$.
Moreover,
\begin{equation}\label{eq:den-identity}
 n^2-j(t)^2=4r_t(n-r_t).
\end{equation}
Admissibility says that every internal $r_t$ differs from both $0$ and $n$.

Fix one admissible path.  For $p=1,\ldots,m+1$, consider the internal vertex
immediately after the $p$-th $+2$ step.  It is internal because the path has
$m+2$ positive steps in total.  Let $d_p\in\{0,1,2,3\}$ be the number of
negative steps that have occurred by then.  The selected vertex is
\begin{equation}\label{eq:rp}
 r_p=2p-d_p.
\end{equation}

For $p=1$, admissibility gives $|r_1|\geq1$.  For $p\geq2$,
\[
 r_p\geq2p-3\geq p-1.
\]
Consequently
\begin{equation}\label{eq:left-product}
 \prod_{p=1}^{m+1}|r_p|\geq\prod_{p=2}^{m+1}(p-1)=m!.
\end{equation}

For the other endpoint factor, set $u=m+1-p$.  From \eqref{eq:rp},
\[
 n-r_p=2u-1+d_p.
\]
When $u\geq1$, this is at least $2u-1\geq u$.  When $u=0$, it is a nonzero
integer by admissibility, hence has absolute value at least one.  Therefore
\begin{equation}\label{eq:right-product}
 \prod_{p=1}^{m+1}|n-r_p|\geq\prod_{u=1}^{m}u=m!.
\end{equation}

There are $m+4$ internal vertices.  The three unselected factors
$r_t(n-r_t)$ are products of nonzero integers and therefore have absolute
value at least one.  Combining \eqref{eq:den-identity},
\eqref{eq:left-product}, and \eqref{eq:right-product} yields the uniform
single-path estimate
\[
 |h_1(x)|\leq\frac{1}{4^{m+4}(m!)^2}.
\]

Every step sequence is determined by the three positions occupied by its
negative steps among $m+5$ positions.  Hence there are at most
$\binom{m+5}{3}$ admissible paths.  Summing the single-path estimate proves
\eqref{eq:main-bound}.
\end{proof}

\begin{corollary}\label{cor:no-C}
One has
\[
 \lim_{m\to\infty}\bigl(m!\,|B_3(2m+1)|\bigr)^{1/m}=0.
\]
In particular, there is no $C>0$ for which
$|B_3(2m+1)|\geq C^m/m!$ holds for all sufficiently large $m$.
\end{corollary}

\begin{proof}
By Theorem~\ref{thm:bound},
\[
 \bigl(m!\,|B_3(2m+1)|\bigr)^{1/m}
 \leq
 \frac{\binom{m+5}{3}^{1/m}}
 {4^{1+4/m}(m!)^{1/m}}\longrightarrow0,
\]
using Stirling's formula.  The final assertion follows immediately.
\end{proof}

\section{Verification and scope}

The argument bounds the absolute sum and is therefore independent of any
cancellation among paths.  The exact rational enumerator in
\texttt{code/verify\_b3.py} checked $m=1,\ldots,12$, including endpoint
admissibility, the signed sum, and the absolute-sum bound.  Run it with
\begin{quote}\small\ttfamily
conda run --no-capture-output -n sandbox\\
python code/verify\_b3.py --max-m 12
\end{quote}
The finite check is not part of the proof.

This packet disproves only the fallback lower estimate in Question~6.  It
does not provide an asymptotic formula for $B_3(2m+1)$, and it does not decide
whether the associated antiperiodic root functions form a basis.  It shows
that inequality~\eqref{eq:source-lower} cannot be used to obtain the proposed
spectral conclusion.

Bounded local-index and arXiv/web searches on 27 August 2026 used the exact
$B_3(2m+1)$ formula, the proposed $C^m/m!$ scale, the paper title and authors,
and the associated two-exponential potential.  They found the source and the
authors' nearby follow-up arXiv:1210.3907, whose stated antiperiodic range does
not include the present $r=1,s=2$ case, but no later treatment of this exact
bound.  Novelty confidence is moderate because the search was bounded and the
obstruction is elementary.

\begin{thebibliography}{1}
\bibitem{DM}
P.~Djakov and B.~Mityagin,
\emph{Combinatorial identities related to eigenfunction decompositions of
Hill operators: Open Questions}, Integral Equations Operator Theory
\textbf{75} (2013), 7--12; arXiv:1210.0835v1,
\url{https://arxiv.org/abs/1210.0835}.
\end{thebibliography}

\end{document}
